\cotefile{157-2}
\scanpage{51}
\mspage{49}
% archive sheet number: 50
Je présume que non -- qu'il peut arriver que
la longueur minimum requise \inserted{$l(n)$} pour un \deleted{\unsure{chemin}}
chemin dans $\underline{\mathrm{Hom}}(K_n,K'_{\unsure{\tau|t}(n)})$ reliant $f_n$ et $g_n$,
tende vers l'$\infty$ si $n\to\infty$. L'hyp que les
$K'_n$ sont des cat. de Thomason y changerait
quelque chose ? (deux objets homotopes dans
$\underline{\mathrm{Hom}}(X,K')$, si $K'$ de Thomason, sont ils reliés
par une flèche ?).

Bien sûr, on peut prétendre que les catégories
$\underline{\mathrm{Ch}}_{\Sigma,n}(M;\unsure{a},\unsure{b})$ ont été définies ``canoniquement'', et
leur limite inductive $\underline{\mathrm{Ch}}_\Sigma(M;a,b)$ aussi.
Mais la définition était assez arbitraire et abracadabrante -- je suis convaincu que dix mathématiciens,
se proposant de développer cette même
idée des catégories de chemins, dans le même
but que moi, arriveraient à dix définitions
différentes ! Pourtant, il n'est pas exclu que
l'on finisse par trouver un système transitif
d'iso. dans $\unsure{\mathrm{Hot}\,\mathrm{hot}|\mathrm{Cat}\,\mathrm{hot}}$, entre les catégories
de chemins \inserted{``infinis''} \deleted{$K_\infty$} (ou de longueur non précisée)
$K_\infty$ \inserted{$(=K_\infty(M;a,b))$} que nous \unsure{aurons|avons} construites, dans le but
